% LaTeX Template for AI Problem Analysis Output
% AMC12B 2024 Problem 13 - Strategic Analysis
% Framework Version: 2.1 (Enhanced for COMC)

\documentclass[12pt]{article}
\usepackage[utf8]{inputenc}
\usepackage{amsmath, amssymb, amsthm}
\usepackage{geometry}
\usepackage{xcolor}
\usepackage[most]{tcolorbox}
\usepackage{enumitem}
\usepackage{fancyhdr}

% Page setup
\geometry{margin=1in}
\setlength{\headheight}{14.5pt}
\pagestyle{fancy}
\fancyhf{}
\rhead{AMC12 Problem Analysis}
\lhead{2024 AMC 12B Problem 13}
\cfoot{\thepage}

% Color definitions
\definecolor{problemconfig}{RGB}{230, 242, 255}    % Light blue
\definecolor{strategic}{RGB}{255, 230, 242}       % Light pink
\definecolor{tactical}{RGB}{242, 255, 230}        % Light green
\definecolor{techniques}{RGB}{255, 242, 230}      % Light yellow
\definecolor{verification}{RGB}{255, 230, 230}    % Light red
\definecolor{solution}{RGB}{230, 255, 255}        % Light cyan
\definecolor{competition}{RGB}{242, 242, 255}     % Light purple
\definecolor{gold}{RGB}{218, 165, 32}

% Box styles
\newtcolorbox{problemconfigbox}{
    colback=problemconfig,
    colframe=blue,
    title=PROBLEM CONFIGURATION ANALYSIS,
    fonttitle=\bfseries,
    arc=3pt,
    boxrule=1pt,
    breakable
}

\newtcolorbox{strategicbox}{
    colback=strategic,
    colframe=purple,
    title=STRATEGIC FRAMEWORK APPLICATION,
    fonttitle=\bfseries,
    arc=3pt,
    boxrule=1pt,
    breakable
}

\newtcolorbox{tacticalbox}{
    colback=tactical,
    colframe=green,
    title=TACTICAL METHODOLOGY SELECTION,
    fonttitle=\bfseries,
    arc=3pt,
    boxrule=1pt,
    breakable
}

\newtcolorbox{techniquesbox}{
    colback=techniques,
    colframe=orange,
    title=CONVERSION TECHNIQUES APPLICATION,
    fonttitle=\bfseries,
    arc=3pt,
    boxrule=1pt,
    breakable
}

\newtcolorbox{verificationbox}{
    colback=verification,
    colframe=red,
    title=VERIFICATION STRATEGY DESIGN,
    fonttitle=\bfseries,
    arc=3pt,
    boxrule=1pt,
    breakable
}

\newtcolorbox{solutionbox}{
    colback=solution,
    colframe=cyan,
    title=SOLUTION ROADMAP,
    fonttitle=\bfseries,
    arc=3pt,
    boxrule=1pt,
    breakable
}

\newtcolorbox{competitionbox}{
    colback=competition,
    colframe=violet,
    title=COMPETITION STRATEGY INTEGRATION,
    fonttitle=\bfseries,
    arc=3pt,
    boxrule=1pt,
    breakable
}

% Special highlight boxes
\newtcolorbox{keyinsightbox}{
    colback=gold!20,
    colframe=gold,
    title=\textbf{KEY INSIGHT},
    fonttitle=\bfseries,
    arc=3pt,
    boxrule=2pt,breakable
}

\newtcolorbox{warningbox}{
    colback=red!10,
    colframe=red!80!black,
    title=\textbf{WARNING},
    fonttitle=\bfseries,
    arc=3pt,
    boxrule=2pt,breakable
}

\newtcolorbox{protipbox}{
    colback=green!10,
    colframe=green!80!black,
    title=\textbf{PRO TIP},
    fonttitle=\bfseries,
    arc=3pt,
    boxrule=2pt
}

\title{\textbf{2024 AMC 12B Problem 13:\\Strategic Analysis}}
\author{Using AMC12/COMC Problem-Solving Framework v2.1}
\date{\today}

\begin{document}

\maketitle

\section*{Problem Statement}

There are real numbers $x, y, h$ and $k$ that satisfy the system of equations

\[x^2 + y^2 - 6x - 8y = h\]
\[x^2 + y^2 - 10x + 4y = k\]

What is the minimum possible value of $h+k$?

\[
\textbf{(A) } -54 \qquad 
\textbf{(B) } -46 \qquad 
\textbf{(C) } -34 \qquad 
\textbf{(D) } -16 \qquad 
\textbf{(E) } 16
\]

\vspace{1em}

\begin{problemconfigbox}
\subsection*{A. Problem Classification}

\begin{itemize}[leftmargin=*]
    \item \textbf{Problem Type:} To Find -- determine the minimum value of an expression
    \item \textbf{Difficulty Band:} Mid-band (Problem 13 of 25)
    \begin{itemize}
        \item Expected time: 3-4 minutes
        \item Requires algebraic manipulation or geometric insight
        \item Tests completing the square and optimization
    \end{itemize}
    \item \textbf{Competition Context:} AMC 12 (75 minutes, 25 problems, multiple choice)
    \item \textbf{Mathematical Domain:} Algebra with Optimization (potential Geometry connection)
    \item \textbf{Source Context:} 2024 AMC 12B -- recent contest problem
\end{itemize}

\subsection*{B. Problem Structure Identification}

\textbf{Given Information:}
\begin{itemize}
    \item Two equations involving $x, y, h, k$
    \item First equation: $x^2 + y^2 - 6x - 8y = h$
    \item Second equation: $x^2 + y^2 - 10x + 4y = k$
    \item All variables are real numbers
    \item We need to find points $(x, y)$ that satisfy both equations for some $(h, k)$
\end{itemize}

\textbf{Constraints:}
\begin{itemize}
    \item $x, y, h, k \in \mathbb{R}$
    \item The system must be consistent (solutions must exist)
    \item We want to minimize $h + k$
\end{itemize}

\textbf{Objective:}
\begin{itemize}
    \item Find $\min(h + k)$
    \item Must match one of five given answer choices
    \item All choices are negative except (E)
\end{itemize}

\textbf{Hidden Assumptions \& Key Observations:}
\begin{itemize}
    \item These equations can be rewritten as circles in standard form
    \item Completing the square will reveal circle centers and radii
    \item Both equations share the $x^2 + y^2$ terms
    \item Adding the equations will eliminate some complexity
    \item The question asks for minimum, suggesting optimization problem
    \item For any $(x, y)$, there exist corresponding $(h, k)$ values
\end{itemize}

\textbf{Answer Choice Leverage:}
\begin{itemize}
    \item Four negative choices, one positive
    \item Magnitude increases from (C) to (A): $-34, -46, -54$
    \item Choice (C) $-34$ is in the middle range
    \item Can verify answer by finding $(x, y)$ that achieves the minimum
\end{itemize}

\subsection*{C. Strategic Orientation}

\textbf{Hypothesis:} 
We need to express $h + k$ as a function of $x$ and $y$, then find the values of $x$ and $y$ that minimize this expression.

\textbf{Conclusion Goal:}
Find the specific value of $\min(h + k)$ and match to answer choices.

\textbf{Notation:}
\begin{itemize}
    \item Let $f(x, y) = h + k$ be the objective function
    \item Circle 1: $(x - 3)^2 + (y - 4)^2 = h + 25$ (after completing square)
    \item Circle 2: $(x - 5)^2 + (y + 2)^2 = k + 29$ (after completing square)
\end{itemize}

\textbf{Key Approaches:}
\begin{enumerate}
    \item \textbf{Direct Algebraic:} Add equations, complete the square, minimize
    \item \textbf{Geometric Interpretation:} View as circles, use geometric properties
\end{enumerate}
\end{problemconfigbox}

\vspace{1em}

\begin{keyinsightbox}
\textbf{Central Insight:} When we add the two equations, we get $h + k = 2x^2 + 2y^2 - 16x - 4y$. This is a quadratic expression in $x$ and $y$ that can be minimized by completing the square. The minimum occurs when both squared terms equal zero, giving us the optimal point $(x, y) = (4, 1)$.
\end{keyinsightbox}

\vspace{1em}

\begin{strategicbox}
\subsection*{A. Psychological Strategies}

\textbf{Mental Toughness:}
\begin{itemize}
    \item System of equations can look intimidating at first
    \item Trust that there's a clean approach (this is mid-band)
    \item Don't overthink -- sometimes the direct algebraic approach is best
\end{itemize}

\textbf{Creativity Cultivation:}
\begin{itemize}
    \item Consider: what happens if we add the two equations?
    \item Recognize the structure: $x^2 + y^2$ terms suggest circles
    \item Think about optimization: minimum of sum of squares occurs at zero
\end{itemize}

\subsection*{B. Investigation Methods}

\textbf{Get Your Hands Dirty -- Try Adding:}
\begin{itemize}
    \item Add equation 1 and equation 2:
    \item $(x^2 + y^2 - 6x - 8y) + (x^2 + y^2 - 10x + 4y) = h + k$
    \item Simplify: $2x^2 + 2y^2 - 16x - 4y = h + k$
    \item Factor: $2(x^2 - 8x) + 2(y^2 - 2y) = h + k$
\end{itemize}

\textbf{Complete the Square:}
\begin{itemize}
    \item For $x$ terms: $x^2 - 8x = (x - 4)^2 - 16$
    \item For $y$ terms: $y^2 - 2y = (y - 1)^2 - 1$
    \item Substitute back
\end{itemize}

\textbf{Recognize Minimum Structure:}
\begin{itemize}
    \item Expression becomes: $h + k = 2(x-4)^2 + 2(y-1)^2 - 34$
    \item Sum of squares is minimized when both squares equal zero
    \item Minimum value: $0 + 0 - 34 = -34$
    \item Occurs at $(x, y) = (4, 1)$
\end{itemize}

\subsection*{C. Alternative Approach: Geometric Interpretation}

\textbf{Recognize Circle Equations:}
\begin{itemize}
    \item Rewrite first equation: $(x - 3)^2 + (y - 4)^2 = h + 25$
    \item Rewrite second equation: $(x - 5)^2 + (y + 2)^2 = k + 29$
    \item These are circles with centers $C_1 = (3, 4)$ and $C_2 = (5, -2)$
    \item Radii: $r_1 = \sqrt{h + 25}$ and $r_2 = \sqrt{k + 29}$ (must be non-negative)
\end{itemize}

\textbf{Geometric Constraint:}
\begin{itemize}
    \item Distance between centers: $d = \sqrt{(5-3)^2 + (-2-4)^2} = \sqrt{4 + 36} = \sqrt{40} = 2\sqrt{10}$
    \item For circles to intersect (share a point): $r_1 + r_2 \geq d$
    \item Therefore: $\sqrt{h + 25} + \sqrt{k + 29} \geq 2\sqrt{10}$
\end{itemize}

\textbf{Optimization via AM-GM:}
This approach is more advanced and less direct for competition use.

\subsection*{D. Argument Construction}

\textbf{Proof Framework (Direct Algebraic):}
\begin{enumerate}
    \item Add the two given equations
    \item Simplify to express $h + k$ in terms of $x$ and $y$
    \item Complete the square for both $x$ and $y$
    \item Recognize that sum of squares $\geq 0$ with equality when squares are zero
    \item Identify the point $(x, y)$ that achieves the minimum
    \item Calculate $\min(h + k)$
\end{enumerate}
\end{strategicbox}

\vspace{1em}

\begin{tacticalbox}
\subsection*{A. Core Mathematical Tactics}

\textbf{Primary Tactic: Completing the Square}
\begin{itemize}
    \item \textbf{Purpose:} Transform quadratic expressions into perfect squares
    \item \textbf{Formula:} $x^2 + bx = (x + \frac{b}{2})^2 - \frac{b^2}{4}$
    \item \textbf{Application:} Makes optimization problems tractable
\end{itemize}

\textbf{Secondary Tactic: Sum of Non-negative Terms}
\begin{itemize}
    \item \textbf{Principle:} $a^2 + b^2 \geq 0$ with equality iff $a = b = 0$
    \item \textbf{Extension:} $c_1(x - a)^2 + c_2(y - b)^2 + k$ has minimum $k$ when $x = a, y = b$
    \item \textbf{Competition Utility:} Immediately identifies optimal point
\end{itemize}

\textbf{Tertiary Tactic: Equation Manipulation}
\begin{itemize}
    \item \textbf{Addition:} Combine equations to create new relationships
    \item \textbf{Subtraction:} Eliminate variables or simplify structure
    \item \textbf{Linear Combination:} Strategic weighting of equations
\end{itemize}

\subsection*{B. Cross-Domain Techniques}

\textbf{Algebra $\leftrightarrow$ Geometry:}
\begin{itemize}
    \item Algebraic equations of form $(x - h)^2 + (y - k)^2 = r^2$ represent circles
    \item Optimization in algebra corresponds to geometric properties (distance, tangency)
    \item Can solve purely algebraically or with geometric insight
\end{itemize}

\textbf{Calculus Connection (Not Needed Here):}
\begin{itemize}
    \item Could use partial derivatives: $\frac{\partial}{\partial x}(h+k) = 0$ and $\frac{\partial}{\partial y}(h+k) = 0$
    \item But completing the square is faster and more elegant for this problem
\end{itemize}

\textbf{Optimization Strategy:}
\begin{itemize}
    \item Transform problem into standard form where minimum is obvious
    \item Leverage structure (sum of squares) rather than brute force calculation
\end{itemize}
\end{tacticalbox}

\vspace{1em}

\begin{protipbox}
\textbf{Competition Hack:} When you see a system of equations where you need to minimize (or maximize) a sum or combination of variables, try adding or subtracting the equations first! This often simplifies the problem dramatically and reveals the structure you need.
\end{protipbox}

\vspace{1em}

\begin{techniquesbox}
\subsection*{A. Recognition Index}

\textbf{Pattern Signature:} System of quadratic equations, minimize sum of parameters

\textbf{Trigger Recognition:}
\begin{itemize}
    \item \textbf{``$x^2 + y^2 + \ldots$''} $\to$ Think circles or completing the square
    \item \textbf{``minimize $h + k$''} $\to$ Look for way to express $h + k$ in terms of $x, y$
    \item \textbf{Two similar equations} $\to$ Try adding/subtracting
    \item \textbf{Quadratic in two variables} $\to$ Complete the square for optimization
\end{itemize}

\subsection*{B. Technique Selection}

\textbf{Approach 1: Direct Algebraic Addition + Completing Square (Recommended)}
\begin{itemize}
    \item \textbf{Why:} Most direct, requires only algebra
    \item \textbf{Setup:} Add equations to get $h + k = 2x^2 + 2y^2 - 16x - 4y$
    \item \textbf{Method:} Complete square for $x$ and $y$, identify minimum
    \item \textbf{Result:} Clean answer $-34$ in under 3 minutes
    \item \textbf{Time:} $\sim$2-3 minutes
\end{itemize}

\textbf{Approach 2: Geometric (Circle) Interpretation}
\begin{itemize}
    \item \textbf{Why:} Provides deeper understanding, connects to geometry
    \item \textbf{Setup:} Rewrite as circles $(x - 3)^2 + (y - 4)^2 = h + 25$, etc.
    \item \textbf{Method:} Use distance between centers and triangle inequality
    \item \textbf{Result:} Same answer but more conceptual steps
    \item \textbf{Time:} $\sim$4-5 minutes
\end{itemize}

\textbf{Approach 3: Calculus (Overkill)}
\begin{itemize}
    \item \textbf{Why:} Would work but unnecessary complexity
    \item \textbf{Method:} Take partial derivatives, set to zero
    \item \textbf{Verdict:} Not recommended for competition setting
\end{itemize}

\subsection*{C. Integration Strategy}

\textbf{Recommended Path for Competition:}
\begin{enumerate}
    \item Read problem, identify need to find $\min(h + k)$ (15 sec)
    \item Recognize both equations have $x^2 + y^2$ terms (10 sec)
    \item Add the two equations (30 sec)
    \item Factor out 2 and complete the square (1 min)
    \item Identify minimum when squares equal zero (15 sec)
    \item Calculate $-34$ and select answer (C) (15 sec)
    \item Quick verification (30 sec)
\end{enumerate}

\subsection*{D. Conflict Resolution}

\textbf{Algebraic vs. Geometric:}
\begin{itemize}
    \item Algebraic approach is faster and more direct
    \item Geometric approach provides insight but takes longer
    \item \textbf{Decision:} Use algebraic unless you immediately see geometric structure
\end{itemize}

\textbf{When to Use Each:}
\begin{itemize}
    \item \textbf{Algebraic:} When time is limited, when you want certainty
    \item \textbf{Geometric:} When you want to understand why, for verification
\end{itemize}
\end{techniquesbox}

\vspace{1em}

\begin{verificationbox}
\subsection*{A. Verification Level Selection}

\textbf{Recommended Level: 3 (Unit Test)}

\textbf{Justification:}
\begin{itemize}
    \item Mid-band problem with straightforward computation
    \item Completing the square is mechanical but can have errors
    \item Answer choices are well-separated
    \item Time budget: approximately 30 seconds for verification
\end{itemize}

\subsection*{B. Verification Methods}

\textbf{Primary Verification: Plug in Optimal Point}
\begin{itemize}
    \item We claim minimum occurs at $(x, y) = (4, 1)$
    \item Calculate $h$ at this point: $h = 16 + 1 - 24 - 8 = -15$
    \item Calculate $k$ at this point: $k = 16 + 1 - 40 + 4 = -19$
    \item Sum: $h + k = -15 + (-19) = -34$ $\checkmark$
\end{itemize}

\textbf{Secondary Verification: Completing Square Check}
\begin{itemize}
    \item Verify: $2(x^2 - 8x + 16) = 2(x - 4)^2$ $\checkmark$
    \item Verify: $2(y^2 - 2y + 1) = 2(y - 1)^2$ $\checkmark$
    \item Constant terms: $-2(16) - 2(1) = -32 - 2 = -34$ $\checkmark$
\end{itemize}

\textbf{Tertiary Verification: Reasonableness Check}
\begin{itemize}
    \item Original equations involve squares and linear terms
    \item After completing square, we get expression with negative constant
    \item Minimum should be negative (sum of squares is non-negative, constant is negative)
    \item $-34$ is reasonable given the structure
    \item Matches answer choice (C) exactly $\checkmark$
\end{itemize}

\subsection*{C. Confidence Calibration}

\textbf{Confidence Level: Very High (98\%)}

\textbf{Reasoning:}
\begin{itemize}
    \item Completing the square is algorithmic and verifiable
    \item Arithmetic is straightforward
    \item Optimal point $(4, 1)$ can be verified by substitution
    \item Multiple verification methods all agree
    \item Answer matches one choice exactly
\end{itemize}

\textbf{Risk Assessment:}
\begin{itemize}
    \item Very low risk of conceptual error (standard technique)
    \item Low risk of arithmetic error (simple fractions and integers)
    \item Very low risk of missing the minimum (squared terms structure is clear)
\end{itemize}
\end{verificationbox}

\vspace{1em}

\begin{warningbox}
\textbf{Common Mistake:} Students might forget to account for the constant terms when completing the square. Remember: $x^2 - 8x = (x - 4)^2 - 16$, not just $(x - 4)^2$. You must subtract the $(\frac{b}{2})^2$ term! This is where arithmetic errors commonly occur.
\end{warningbox}

\vspace{1em}

\begin{solutionbox}
\subsection*{Solution Roadmap -- Primary Approach (Direct Algebraic)}

\textbf{Step 1: Add the Two Equations (30 seconds)}

Given:
\begin{align*}
x^2 + y^2 - 6x - 8y &= h \quad \text{(Equation 1)}\\
x^2 + y^2 - 10x + 4y &= k \quad \text{(Equation 2)}
\end{align*}

Add them:
\begin{align*}
(x^2 + y^2 - 6x - 8y) + (x^2 + y^2 - 10x + 4y) &= h + k\\
2x^2 + 2y^2 - 16x - 4y &= h + k
\end{align*}

\textbf{Step 2: Factor Out Common Terms (15 seconds)}
\begin{align*}
h + k &= 2x^2 + 2y^2 - 16x - 4y\\
&= 2(x^2 - 8x) + 2(y^2 - 2y)
\end{align*}

\textbf{Step 3: Complete the Square for $x$ (30 seconds)}

For $x^2 - 8x$:
\begin{itemize}
    \item Take half of $-8$: $\frac{-8}{2} = -4$
    \item Square it: $(-4)^2 = 16$
    \item Rewrite: $x^2 - 8x = (x - 4)^2 - 16$
\end{itemize}

Therefore:
\[2(x^2 - 8x) = 2[(x - 4)^2 - 16] = 2(x - 4)^2 - 32\]

\textbf{Step 4: Complete the Square for $y$ (30 seconds)}

For $y^2 - 2y$:
\begin{itemize}
    \item Take half of $-2$: $\frac{-2}{2} = -1$
    \item Square it: $(-1)^2 = 1$
    \item Rewrite: $y^2 - 2y = (y - 1)^2 - 1$
\end{itemize}

Therefore:
\[2(y^2 - 2y) = 2[(y - 1)^2 - 1] = 2(y - 1)^2 - 2\]

\textbf{Step 5: Combine Results (30 seconds)}
\begin{align*}
h + k &= 2(x - 4)^2 - 32 + 2(y - 1)^2 - 2\\
&= 2(x - 4)^2 + 2(y - 1)^2 - 34
\end{align*}

\textbf{Step 6: Identify Minimum (15 seconds)}

The expression $h + k = 2(x - 4)^2 + 2(y - 1)^2 - 34$ is a sum of:
\begin{itemize}
    \item $2(x - 4)^2 \geq 0$ (equality when $x = 4$)
    \item $2(y - 1)^2 \geq 0$ (equality when $y = 1$)
    \item Constant term: $-34$
\end{itemize}

Minimum value occurs when both squared terms equal zero:
\begin{itemize}
    \item $(x - 4)^2 = 0 \Rightarrow x = 4$
    \item $(y - 1)^2 = 0 \Rightarrow y = 1$
\end{itemize}

At $(x, y) = (4, 1)$:
\[\min(h + k) = 2(0) + 2(0) - 34 = -34\]

\textbf{Step 7: Match to Answer Choice (5 seconds)}
\begin{itemize}
    \item Our result: $-34$
    \item This is answer choice \textbf{(C)}
\end{itemize}

\subsection*{Time Checkpoint}

\textbf{Total Time: Approximately 2.5-3 minutes}
\begin{itemize}
    \item Adding equations: 30 sec
    \item Factoring: 15 sec
    \item Completing square for $x$: 30 sec
    \item Completing square for $y$: 30 sec
    \item Combining: 30 sec
    \item Finding minimum: 15 sec
    \item Verification: 30 sec
\end{itemize}

This is faster than expected for a Problem 13, indicating an elegant solution approach!

\subsection*{Verification of Answer}

Let's verify by substituting $(x, y) = (4, 1)$ back into the original equations:

\textbf{For $h$:}
\begin{align*}
h &= x^2 + y^2 - 6x - 8y\\
&= 16 + 1 - 24 - 8\\
&= -15
\end{align*}

\textbf{For $k$:}
\begin{align*}
k &= x^2 + y^2 - 10x + 4y\\
&= 16 + 1 - 40 + 4\\
&= -19
\end{align*}

\textbf{Sum:}
\[h + k = -15 + (-19) = -34 \quad \checkmark\]

Perfect match!

\subsection*{Alternative Approach: Geometric Interpretation}

\textbf{Rewrite as Circles:}

Complete the square in each equation separately:

\textbf{Equation 1:}
\begin{align*}
x^2 + y^2 - 6x - 8y &= h\\
(x^2 - 6x) + (y^2 - 8y) &= h\\
(x^2 - 6x + 9) + (y^2 - 8y + 16) &= h + 9 + 16\\
(x - 3)^2 + (y - 4)^2 &= h + 25
\end{align*}

This represents a circle with center $C_1 = (3, 4)$ and radius $r_1 = \sqrt{h + 25}$ (assuming $h \geq -25$).

\textbf{Equation 2:}
\begin{align*}
x^2 + y^2 - 10x + 4y &= k\\
(x^2 - 10x) + (y^2 + 4y) &= k\\
(x^2 - 10x + 25) + (y^2 + 4y + 4) &= k + 25 + 4\\
(x - 5)^2 + (y + 2)^2 &= k + 29
\end{align*}

This represents a circle with center $C_2 = (5, -2)$ and radius $r_2 = \sqrt{k + 29}$ (assuming $k \geq -29$).

\textbf{Distance Between Centers:}
\begin{align*}
d &= \sqrt{(5 - 3)^2 + (-2 - 4)^2}\\
&= \sqrt{4 + 36}\\
&= \sqrt{40}\\
&= 2\sqrt{10}
\end{align*}

\textbf{Geometric Constraint:}

For the system to have a solution (i.e., for there to exist a point $(x, y)$ on both circles), the circles must intersect. By the triangle inequality:
\[r_1 + r_2 \geq d\]
\[\sqrt{h + 25} + \sqrt{k + 29} \geq 2\sqrt{10}\]

The minimum of $h + k$ occurs when this inequality is tight (circles are tangent).

\textbf{Optimization:} Using techniques like AM-GM or Lagrange multipliers (beyond typical competition scope), one can show that:
\[h + k \geq -34\]

with equality at $(x, y) = (4, 1)$.

\textbf{Finding the Tangent Point:}

The point $(4, 1)$ lies on the line segment connecting $C_1 = (3, 4)$ and $C_2 = (5, -2)$:
\begin{itemize}
    \item Parametric line: $(x, y) = (3, 4) + t(2, -6)$ for $t \in [0, 1]$
    \item At $t = \frac{1}{2}$: $(x, y) = (3, 4) + \frac{1}{2}(2, -6) = (4, 1)$ $\checkmark$
\end{itemize}

This confirms that $(4, 1)$ is the midpoint of the line segment joining the two centers, which makes geometric sense for the tangency condition.

\textbf{Complexity Assessment:} The geometric approach provides beautiful insight but requires more steps than the direct algebraic method. It's excellent for understanding the problem deeply but less efficient for competition.
\end{solutionbox}

\vspace{1em}

\begin{competitionbox}
\subsection*{A. Time Management}

\textbf{Time Allocation for Problem 13:}
\begin{itemize}
    \item Target time: 2.5-3 minutes
    \item Maximum time: 4 minutes
    \item This problem is faster than its position suggests
\end{itemize}

\textbf{Decision Point:}
\begin{itemize}
    \item If you see the "add equations" approach immediately: proceed (2.5 min)
    \item If you think about geometry: it works but takes longer (4-5 min)
    \item If stuck after 1 minute: try adding the equations as a heuristic
\end{itemize}

\subsection*{B. Answer Format Requirements}

\textbf{AMC 12 Multiple Choice:}
\begin{itemize}
    \item Select one answer: \textbf{(C)} $-34$
    \item No partial credit
    \item Guessing penalty: None (score 0 for wrong/blank)
\end{itemize}

\subsection*{C. Strategic Context}

\textbf{Problem 13 Position Strategy:}
\begin{itemize}
    \item This is early-to-mid mid-band territory
    \item Most students with algebra proficiency should attempt
    \item Correct approach yields answer very quickly (2-3 minutes)
    \item Worth 6 points (standard for all AMC problems)
    \item Tests fundamental skill: completing the square
\end{itemize}

\textbf{Skills Tested:}
\begin{itemize}
    \item Algebraic manipulation (adding equations)
    \item Completing the square (fundamental technique)
    \item Optimization (recognizing minimum of sum of squares)
    \item Verification (checking answer makes sense)
\end{itemize}

\subsection*{D. Risk Assessment}

\textbf{Risk Level: Low}

\textbf{Factors:}
\begin{itemize}
    \item \textbf{Low Risk:} Standard technique with mechanical steps
    \item \textbf{Pitfall:} Arithmetic errors in completing the square
    \item \textbf{Mitigation:} Double-check constant terms
\end{itemize}

\textbf{When to Skip:}
\begin{itemize}
    \item If completing the square is unfamiliar (review needed)
    \item If already severely behind on time
    \item Very rare case for Problem 13 position
\end{itemize}

\textbf{When to Commit:}
\begin{itemize}
    \item If you're comfortable with completing the square
    \item If you have 3+ minutes available
    \item If you need a confidence boost (this problem is very doable)
\end{itemize}

\subsection*{E. Pattern Recognition for Similar Problems}

\textbf{This Problem Type Appears When:}
\begin{itemize}
    \item System of equations with quadratic terms
    \item Asked to optimize (minimize or maximize) a combination
    \item Completing the square is the key technique
\end{itemize}

\textbf{General Strategy:}
\begin{enumerate}
    \item Try adding/subtracting equations
    \item Look for common terms to eliminate
    \item Complete the square to find optimal points
    \item Use non-negativity of squared terms
\end{enumerate}
\end{competitionbox}

\vspace{1em}

\begin{protipbox}
\textbf{Post-Exam Learning:} This problem exemplifies a key AMC strategy: when you have a system of similar equations and need to optimize a combination of variables, manipulating the equations (adding, subtracting, etc.) often reveals structure that makes the problem trivial. Always ask: ``What happens if I add/subtract these?''
\end{protipbox}

\newpage

\section*{Complete Solution Summary}

\subsection*{Key Steps:}

\begin{enumerate}
    \item \textbf{Add the Equations:} 
    
    Sum the two given equations to express $h + k$:
    \[h + k = 2x^2 + 2y^2 - 16x - 4y\]
    
    \item \textbf{Factor and Prepare for Completing Square:}
    \[h + k = 2(x^2 - 8x) + 2(y^2 - 2y)\]
    
    \item \textbf{Complete the Square for $x$:}
    \[2(x^2 - 8x) = 2(x - 4)^2 - 32\]
    
    \item \textbf{Complete the Square for $y$:}
    \[2(y^2 - 2y) = 2(y - 1)^2 - 2\]
    
    \item \textbf{Combine:}
    \[h + k = 2(x - 4)^2 + 2(y - 1)^2 - 34\]
    
    \item \textbf{Identify Minimum:}
    
    The minimum occurs when both squared terms equal zero:
    \begin{itemize}
        \item $x = 4$
        \item $y = 1$
    \end{itemize}
    
    \item \textbf{Calculate Minimum Value:}
    \[\min(h + k) = 0 + 0 - 34 = -34\]
\end{enumerate}

\subsection*{Final Answer:}

\begin{center}
\fbox{\Large \textbf{(C) } $\displaystyle -34$}
\end{center}

\vspace{1em}

\subsection*{Verification:}
\begin{itemize}
    \item At $(x, y) = (4, 1)$: $h = 16 + 1 - 24 - 8 = -15$ $\checkmark$
    \item At $(x, y) = (4, 1)$: $k = 16 + 1 - 40 + 4 = -19$ $\checkmark$
    \item Sum: $h + k = -15 + (-19) = -34$ $\checkmark$
    \item Completing square arithmetic verified $\checkmark$
    \item Matches answer choice (C) exactly $\checkmark$
\end{itemize}

\subsection*{Key Takeaways:}

\begin{enumerate}
    \item \textbf{Equation Manipulation:} When dealing with systems of equations and optimization, try combining equations (adding, subtracting) to reveal structure.
    
    \item \textbf{Completing the Square:} This is a fundamental technique for optimization problems. Remember to account for constant terms: $x^2 + bx = (x + \frac{b}{2})^2 - \frac{b^2}{4}$.
    
    \item \textbf{Sum of Squares Optimization:} Expressions of the form $a(x - h)^2 + b(y - k)^2 + c$ (with $a, b > 0$) have minimum value $c$ at $(x, y) = (h, k)$.
    
    \item \textbf{Geometric Interpretation:} While the algebraic approach is faster, recognizing that the equations represent circles provides geometric insight and an alternative verification method.
    
    \item \textbf{Verification by Substitution:} Always check your answer by substituting the optimal point back into the original equations.
    
    \item \textbf{Speed vs. Insight:} Sometimes the fastest method (algebraic) and the most insightful method (geometric) differ. In competition, prioritize speed; in study, explore both.
\end{enumerate}

\vspace{2em}

\hrule

\vspace{1em}

\begin{center}
\textit{Analysis completed using AMC12/COMC Strategic Problem-Solving Framework v2.1}

\textit{Framework emphasizes: Algebraic Manipulation • Optimization Techniques • Multiple Perspectives}
\end{center}

\end{document}